% titel: Blatt 3, A2: Binomialverteilung

% Teil I: Werkzeuge, kein Lösungsweg.

Zwei konkrete Wahrscheinlichkeiten mit der Binomialverteilung bestimmen:
verspätete Zugfahrten und herausspringende Fische.

\begin{werkzeugkasten}[Binomialverteilung]
Bei $n$ unabhängigen Versuchen mit Erfolgswahrscheinlichkeit $p$ ist die
Wahrscheinlichkeit für genau $k$ Erfolge
\begin{equation}
  \mathrm P(k)=\binom{n}{k}p^k(1-p)^{n-k}.
\end{equation}
\begin{itemize}
  \item ``Mindestens $4$ von $5$ verspätet'' (Erfolg $=$ Verspätung,
        $p=\tfrac15$): $\mathrm P(4)+\mathrm P(5)$.
  \item ``Mindestens einer von $n$'': Komplement
        $1-\mathrm P(\text{keiner})=1-(1-p)^n$.
\end{itemize}
\end{werkzeugkasten}

\begin{achtung}[Darauf achten]
\begin{itemize}
  \item ``Mindestens $4$ von $5$'' heißt $4$ \emph{oder} $5$ -- beide Fälle
        addieren ($k=4$ und $k=5$).
  \item ``Mindestens ein Fisch springt hinaus'' bequemer über das
        Gegenereignis ``keiner springt hinaus'' $=(1-p)^n$.
  \item $p$ und $\binom{n}{k}$ korrekt einsetzen: $\mathrm P(5)=\binom55 p^5 (1-p)^0=p^5$.
\end{itemize}
\end{achtung}
